\subsection{Le modèle des facteurs de vue d'ordre 1}

Les facteurs de vue sont introduits dans le but de simplifier l'écriture de l'irradiance (Équation~\ref{def:irra}), en faisant apparaître des facteurs purement géométriques qui décrivent le couplage radiatif entre surfaces. Cette simplification repose sur une hypothèse supplémentaire, propre au modèle dit \emph{d'ordre 1} : l'\textbf{uniformité de la radiosité sur chaque surface}. Sous cette hypothèse, la radiosité de la surface $i$ ne dépend plus de la position $\vec{r}_i$ mais se réduit à une valeur moyenne $\overline{J_i^1}$~[W/m²]. Elle permet de faire apparaître un terme géométrique $f_{ij}$ dans l'expression de l'irradiance, tel que pour $\forall i$ :

\begin{equation}
    G_i(\vec{r_i}) = \sum_{j = 1}^{n} \overline{J_j}^1 f_{ij}(\vec{r_i}) \quad \text{avec } f_{ij}(\vec{r_i}) =\frac{1}{\pi} \int_{A_i} d^2(\vec{r_j})\bigg[\frac{cos(\theta_{ij})cos(\theta_{ji})}{||\vec{r_i}-\vec{r_j}||^2} \chi{(\vec{r_i},\vec{r_j})}\bigg] 
\end{equation}

L'intégration sur la surface $i$ permet d'obtenir une approximation à l'ordre 1 de la fraction du rayonnement quittant la surface $i$ interceptée par la surface $j$. Cette quantité est appelée facteur de vue (\textit{les termes de facteur de forme ou facteur de configuration sont aussi trouvés dans la littérature}).

\begin{equation}
    \label{def:fv}
    F_{ij} = \frac{1}{A_i} \int_{A_i} \int_{A_j} d^2(\vec{r_i}) d^2(\vec{r_j})\bigg[\frac{cos(\theta_{ij})cos(\theta_{ji})}{\pi ||\vec{r_i}-\vec{r_j}||^2} \chi(\vec{r_i},\vec{r_j})\bigg] 
\end{equation}


L'hypothèse de radiosité uniforme revient ainsi à ramener le problème radiatif continu à un problème d'éclairement mutuel entre un nombre fini de surfaces. La connaissance des facteurs de vue permet alors de reformuler l'équation \ref{def:radiosite} pour exprimer les radiosités en fonction de la seule variable température :

\begin{equation}
    \overline{J_i^1} = \epsilon_i \sigma \overline{T}^4_i  + \rho_i \sum_{j=1}^{n} \overline{J_j^1} F_{ij}
    \label{def:radio}
\end{equation}


On remarque par ailleurs que les facteurs de vue respectent des propriétés physiques fondamentales, à savoir :

\begin{equation}    
    \left\lbrace
    \begin{array}{r c c c  }
      \forall i ,\forall j, A_i F_{ij} & = & A_j F_{ji} & \text{(relation de réciprocité)}\\
      \forall i , \sum_{j = 1}^n F_{ij} & = & 1 & \text{(relation de fermeture)}\\ 
      \forall i ,\forall j, F_{ij} & \geq & 0 & \text{(positivité)}
    \end{array}\right.
    \label{eq:contraintes}
\end{equation}

Enfin, l'Équation \ref{eq:qr} qui définit le flux radiatif devient :

\begin{equation}
    \forall \vec{r}_i \in \text{Surface}_i,\quad q_i(\vec{r}_i) =  \overline{J}_i^1  -  \sum_{j=1}^{n} \overline{J_j^1} F_{ij} \quad = \quad \epsilon_i \sigma \overline{T}^4_i  - \epsilon_i \sum_{j=1}^{n} \overline{J_j^1} F_{ij}
    \label{eq:nice_qr}
\end{equation}